% S9 manuscript — JEDM official template conversion (paper-jedm.tex)
% Template: jedm.cls from the official JEDM LaTeX package 2024_08_02_jedm-latex.zip,
% downloaded from the journal's Information for Authors page
% (https://jedm.educationaldatamining.org/index.php/JEDM/information/authors).
% Converted from paper.tex (article class with hand-rolled JEDM-spec formatting).
% FORMAT CONVERSION ONLY — every word of content, number, citation key,
% figure/table, and the manual bibliography is preserved verbatim from paper.tex.
% Anonymized for double-blind review (empty author block).
\documentclass[notitlepage,12pt]{jedm}

\usepackage{mathptmx}              % Times Roman body + math
\usepackage[scaled=0.92]{helvet}   % Helvetica for title/headings (JEDM spec)
\usepackage{graphicx}
\usepackage{booktabs}
\usepackage{amsmath,amssymb}
% jedm.cls pre-allocates \bibhang; release it so natbib can allocate its own
\let\bibhang\relax
\usepackage[round]{natbib}
\usepackage[hidelinks]{hyperref}
\usepackage{microtype}

\title{No Learning Under This Protocol:\\A Preregistered Validity Audit of a Downsized Frozen-Student Replication of Sherpa's Simulator}
\author{}   % anonymized; double-blind submission
\date{}
\begin{document}
\maketitle

\begin{abstract}
\noindent Large language models are increasingly used as simulated students:
synthetic learners that should mirror real students, used to train
tutors and study pedagogy. We report a preregistered, independent validity audit of the student
simulator (Sherpa), which conditions a weight-frozen (weights held fixed) language model on six
preference archetypes via gating mechanisms (preference filters on teacher
messages). Replicating only the student
side, we generated 300 episodes (six archetypes, 50 trajectories each, five
tutoring turns) under a fixed-strategy teacher and compared behavioral
statistics against five thousand real student sequences from the
ASSISTments 2009--2010 skill-builder dataset. The simulated students
do not learn at all: a correctness probe (re-asking the student to solve the
problem after tutoring) is frozen behaviorally (constant
across turns) at its baseline value in 297
of 300 trajectories, and the 24 preregistered two-sample tests are
degenerate by construction---the simulated metrics collapse to point masses
while real students show the classic learning curve, so their p-values add
nothing beyond the descriptive mismatch. No behavioral metric
separates the six archetypes. Gate decomposition shows nearly
half of adaptive-gate rejections are fail-closed judge parse failures
rather than preference judgments. A post-hoc positive control---injecting
the ground-truth boxed answer (the convention of
enclosing the final answer in a \texttt{\textbackslash boxed\{\}} tag)
into the probe's history in 162 of 300 episodes (four of six archetypes)---moves
zero of the wrong-at-baseline probes to correct. This partial, post-hoc
control cannot separate instruction-following from incapacity: the
``solve from scratch'' probe instruction may simply tell the student to
ignore history. The freeze therefore cannot be attributed
to preference gating specifically, and our conclusions are scoped to this
downsized protocol.
{\parindent0pt
\textbf{Keywords:} large language models, preregistration, simulated students, validity audit
}
\end{abstract}

\section{Introduction}

Simulated students---language-model agents designed to behave like human
learners---are becoming infrastructure. They train tutoring systems
\citep{xu2026sherpa}, pilot instructional interventions, and stand in for
human subjects in education research. Their scientific value rests entirely
on a validity claim: that the simulator's behavior mirrors real student
behavior closely enough for conclusions to transfer. That claim is rarely
audited independently. The stakes are concrete: a tutoring policy trained against
a simulator that cannot learn optimizes pedagogy signals that would never
transfer to real students---compute spent training on an instrument that
measures nothing.

Sherpa \citep{xu2026sherpa} proposes archetype-conditioned simulated
students: a frozen language model whose teaching interactions are filtered
through preference gates (enforced in the paper's pipeline by a 27B
language-model judge around an 8B teacher), producing six ``distinct'' student archetypes (for
example, attempt-diagnosis and subgoal-decomposition). The paper's evaluation
targets the \emph{teacher} it trains---student accuracy gains, MathTutorBench
pedagogy scores, and a human preference study---while the student simulator's
own validity as a learner goes untested---most basically, whether the
simulated students learn like students. Real students exhibit
one of the most robust phenomena in cognitive science---the power law of
practice, in which error rates decay with opportunity \citep{newell1981}. A
simulator that cannot reproduce it is not a student model in any behaviorally
meaningful sense. Our probes re-attempt the \emph{same} problem after
tutoring, so the test here is the minimal within-problem version---whether
instruction leaves any trace on the next attempt (retention), not transfer
across problems; a simulator failing even this bar fails the stronger claim
a fortiori.

We preregistered an independent audit of Sherpa's student simulator before
generating a single trajectory (preregistration locked 2026-10-08; one logged
deviation). Replicating only the student side---frozen Qwen3-1.7B with
verbatim gate prompts from the released repository---we generated 300
episodes (6 archetypes $\times$ 50 trajectories $\times$ 5 tutoring turns)
under a fixed-strategy teacher built from the paper's own teaching templates,
and compared behavioral statistics against 5{,}000 real (student, skill)
sequences from the public ASSISTments 2009--2010 skill-builder dataset
\citep{feng2009}.

Our findings are uniformly negative:

\begin{enumerate}
\item \textbf{Zero learning.} Probe correctness is frozen at its unaided
baseline in 297 of 300 trajectories. Mean correctness moves 0.347
$\rightarrow$ 0.343 across five tutoring turns. Exactly one trajectory
improved---with zero accepted tutoring turns (probe noise; see
\S\ref{sec:res-probe}). Two got
\emph{worse} after accepted tutoring.
\item \textbf{Degenerate distributional mismatch.} All 24 preregistered
two-sample Kolmogorov--Smirnov tests (4 metrics $\times$ 6 archetypes,
Holm-adjusted) reject the null. But the ``difference'' is degenerate:
the simulated distributions collapse to point masses---slopes $0.000$
$\pm$ 0.000, error persistence $0.833$ $\pm$ 0.000, hint
responsiveness $0.000$ $\pm$ 0.000 across four archetypes, with only the three
mixed trajectories giving the other two any spread---while real students show the
classic learning curve (0.499 $\rightarrow$ 0.688, mean slope $-0.35$).
\item \textbf{Indistinguishable archetypes.} On every behavioral metric, the
six archetypes are identical. The only differences are gate acceptance rates
(20--46\%) and surface reply lengths.
\item \textbf{Gate dysfunction, decomposed.} Of 1{,}500 teacher turns, 17.2\%
are rejected by the guidance gate, 24.7\% by adaptive-gate \emph{parse
failures} (the 1.7B judge cannot emit valid tagged XML after three retries;
fail-closed), and 26.3\% by reasoned preference rejections; 31.9\% are
accepted. Nearly half the adaptive rejection mass is judge incompetence, not
preference enforcement.
\item \textbf{No dose--response.} The correlation between accepted tutoring
turns per episode and probe improvement is $r=-0.04$ (95\% CI $[-0.15,
0.08]$, $n=300$). Five accepted turns move the probe about as much as
zero---and in 23 episodes where the teacher's boxed final answer was
\emph{accepted} into the student history, zero probes improved. Post-hoc
checking revealed every one of the 45 leaked boxed answers in the 32
accepted messages was numerically wrong (0/45 correct), so this null
indicts the 1.7B teacher's quality rather than testing whether tutoring
content can move probes (\S\ref{sec:res-boxed}).
\end{enumerate}

This is a negative-validity audit, framed as such: the contribution is not a
better simulator but a preregistered demonstration
that this simulator design, as implemented, exhibits no
student-like learning dynamics whatsoever. The gate judge and teacher are
downsized to 1.7B (a CPU constraint) versus the paper's 27B judge and 8B
teacher, which explains part of the gate dysfunction---but it cannot explain
the frozen probes, because even injected true answers move nothing
(\S\ref{sec:res-control}).

\section{Related work}
\label{sec:related}

\paragraph{Language-model simulated students.} The use of language models as
stand-ins for human learners has grown rapidly, but validity evidence has
lagged behind adoption. \citet{scarlatos2026} formally define the
student-simulation task and benchmark simulation methods on linguistic,
behavioral, and cognitive metrics, finding that prompting-based simulators
perform poorly; \citet{yuan2026} identify a ``competence paradox''---broadly
capable language models asked to emulate partially knowledgeable learners
produce unrealistic error patterns and learning dynamics---and argue for
epistemic fidelity over surface realism. Sherpa \citep{xu2026sherpa} is
among the most elaborate recent systems: six preference archetypes enforced
by language-model judges (a guidance gate plus an adaptive gate). Its
evaluation targets the \emph{teacher}---student accuracy gains,
MathTutorBench pedagogy scores, and a human pairwise-preference study---while
the student simulator's own behavioral validity \emph{as a learner} is
assumed, not tested. Our audit is, to our knowledge, the first
\emph{independent, preregistered} test of whether a recent
simulated-student system reproduces real student \emph{learning dynamics}
against real student data: prior validity work targets static response
patterns, not learning over time. We test not whether the
simulator is internally coherent, but whether it learns like a student at
all.

Validation efforts elsewhere cluster around response patterns rather than
learning dynamics. \citet{lu2024} build knowledge-component-grounded
``generative students'' whose simulated responses correlate highly with real
students' item difficulties---a positive validity result, but for static
response patterns, not learning over time. \citet{srivatsa2025} place eleven
language models on the same IRT ability scale as real student populations
and find systematic misalignment: strong models overshoot every grade level.
\citet{an2026} identify the mechanism: prompt-based simulators answer at
their built-in capability level (96.8--100\% accuracy) regardless of the
assigned mastery profile. \citet{scarlatos2025} bring knowledge tracing into
tutor--student dialogues with language models. These simulators are deployed
as training environments for tutors \citep{xu2026sherpa}, as proxy students
for piloting assessments \citep{srivatsa2025}, and as stand-ins for real
respondents in item evaluation \citep{lu2024}---all roles that presuppose
student-like behavior. None test whether a simulator's behavior
\emph{changes} like a student's under instruction: the learning-dynamics gap
our audit targets.

\paragraph{Student modeling and learning curves.} The power law of
practice---error rate decaying as a power function of opportunities---is
among the most replicated findings in cognitive science \citep{newell1981}
and underpins knowledge-tracing models from Bayesian knowledge tracing
\citep{corbett1994} onward. The ASSISTments datasets \citep{feng2009} have
anchored decades of student-modeling research. Any claim to simulate
students inherits this literature as its validity baseline.

\paragraph{Audits and replication.} Independent audits of high-profile
machine-learning artifacts have repeatedly surfaced validity gaps, and
preregistered replication is the mechanism by which those gaps become
credible rather than anecdotal \citep{nosek2018}. We adopt that mechanism
here.

\section{Preregistration}
\label{sec:prereg}

The full preregistration (locked 2026-10-08, before any study trajectory was
generated) is archived with the reproduction package. In brief:

\paragraph{Study question.} Do Sherpa's archetype-conditioned simulated
students reproduce real student learning dynamics? In a validity audit
the roles reverse: the hypothesis of interest is mismatch, and the
preregistered tests are a decision-rule formality---the evidential weight
lies in the descriptive degeneracy. H1: simulated students
differ measurably from real students on at least one of (M1) learning-curve
shape, (M2) error persistence, (M3) hint responsiveness. H0: no significant
distributional difference on any preregistered metric (24 two-sample KS
tests, Holm--Bonferroni, $\alpha=0.05$).

\paragraph{Protocol.} Six preference archetypes---attempt-diagnosis (Att.),
subgoal-decomposition (Sub.), contrastive-comparison (Ctr.),
causal-justification (Cau.), step-demonstration (Stp.),
independent-verification (Ver.)---implemented with frozen Qwen3-1.7B, a
``you are Sam, a student'' system message, and the verbatim adaptive-gate
preference text from the paper's Appendix D.1.2). Per episode: an unaided
baseline probe ($c_0$), then five turns of fixed-strategy teacher instruction
(the paper's Appendix D.5.1 templates, strategy matched to the archetype's
preference) filtered through the verbatim guidance gate and then the
adaptive gate (up to three parse retries, fail-closed, per the released
repository), then a greedy probe fork ($c_t$) graded by the repository's
two-stage scorer. Fifty problems per archetype, disjoint, drawn from the
paper's shipped filtered MATH test split (528 problems). Real-data
comparison: 5{,}000 seeded (user, skill) sequences from the ASSISTments
2009--2010 skill-builder data (first five opportunities, main problems
only).

\paragraph{Metrics.} M1: per-trajectory power-law exponent from
$\log(e_t)=\alpha+b\log(t)$ on smoothed error rates (constant probe vectors
yield slope exactly 0 by construction)---it measures whether the student
improves with practice, the learning-curve signature. M2a:
Laplace-smoothed $P(\text{err}_{t+1}\mid\text{err}_t)$---whether mistakes
persist or get corrected (for constant probe vectors this equals
$(4+1)/(4+2)=0.833$ by the smoothing, a definitional constant rather than
an estimated persistence). M2b: maximum error streak---the depth of
uncorrected failure runs. M3: for simulated trajectories, mean($c_t\mid$guidance accepted) minus
mean($c_t\mid$guidance rejected); for real trajectories, next-opportunity
correctness after
a hint versus after no hint---whether instruction moves subsequent
performance. Full definitions are in the preregistration.

\paragraph{One deviation.} On 2026-10-09, the context window was raised from
2048 to 4096 tokens after a deterministic overflow at episode 31 crashed
every relaunch. Context size does not change model logits for prompts that
fit in 2048 tokens (verified post hoc from the banked artifacts: no single
model call in the 31 banked episodes exceeded 1{,}470 prompt tokens), so the 31 banked episodes remain comparable; the
alternative (dropping the episode) would have broken the 300-episode design.
The post-hoc answer-injection control (\S\ref{sec:res-control}) is likewise
logged in the deviation log as a post-hoc addition: it was stopped after
162 of 300 episodes (2-CPU compute contention at roughly five minutes per
episode; episodes processed in file order, so the last two archetypes never
ran).

\paragraph{Stated limitations.} Competition-mathematics versus K-12 domain
mismatch; downsized 1.7B judge and teacher (CPU constraint); five turns
versus the paper's ten; Q4\_K\_M quantization (taken as a stated
assumption). See \S\ref{sec:limitations} for the full discussion.

\paragraph{Preregistration boundary.} The preregistration covers the 24 KS
tests and the metrics in \S\ref{sec:res-learn}--\S\ref{sec:res-arch}. The gate
decomposition (\S\ref{sec:res-gates}), boxed-answer and dose--response
analyses (\S\ref{sec:res-boxed}), the probe-stochasticity check
(\S\ref{sec:res-probe}), and the answer-injection positive control
(\S\ref{sec:res-control}) are exploratory, post-hoc analyses conducted after
the frozen-probe result was observed; they are labeled as such and carry no
preregistered decision rule.

\section{Methods}
\label{sec:methods}

\begin{center}
\begin{minipage}{0.94\columnwidth}
\hrule\vspace{0.6em}
\noindent\textbf{Protocol at a glance.} \textbf{Probe:} a greedy
(temperature-0.0) re-attempt of the \emph{same} problem by the frozen
student, run on a fork of the accepted-turn history; the unaided baseline
probe ($c_0$) sees no tutoring, each post-tutoring probe ($c_1$--$c_5$)
sees the full accepted history. \textbf{Gate:} the two-stage filter---the
guidance gate, then the archetype's adaptive gate---deciding whether a
teacher turn enters the student's visible history; rejected turns are
hidden and replaced by a scripted complaint. \textbf{Fork:} every probe
branches off the current accepted history instead of continuing the
dialogue, so probe correctness measures whether the tutoring history
taught anything.
\vspace{0.6em}\hrule
\end{minipage}
\end{center}

\subsection{Simulation}

All generation used llama-cpp-python on CPU (Qwen3-1.7B Q4\_K\_M GGUF).
Generation parameters were locked in the preregistration: teacher sampling
at temperature 0.7, gates at 0.0, student replies at 0.7, probes greedy
(temperature 0.0), intended as deterministic state read-outs---an assumption
that is violated: greedy probes are generation-stochastic, though their
graded verdicts are stable (see \S\ref{sec:res-probe})---and a temperature-0.0
answer-judge fallback. Each probe is a fork of the current accepted-turn
history: the baseline probe ($c_0$) sees only the system prompt, while each
post-tutoring probe ($c_t$) re-attempts the task with the full accepted
tutoring history in context (rejected turns excluded), so tutoring that
teaches would move the probe. Thinking mode was disabled in the paper. The master
seed was 20261008 with per-episode seeds derived by hashing. Rejected turns
never enter the student-visible history (the repository's hidden-from-student
turn semantics); the student instead emits the paper's scripted complaint.
Every generation is logged with seeds (more than 7{,}800 generations in
total).

\paragraph{Filter fidelity.} The shipped filtered split is described as
retaining problems the student fails unaided, yet our greedy baseline
probes answer 34.7\% correctly---which is not a filter failure. Sherpa's
own Table~1 reports the student's overall unaided accuracy on these
problems as 36\%, and our 34.7\% $\pm$ 2.7\% reproduces it ($z=-0.49$
against 36\%). The paper's unaided estimate comes from 8 repeated samples
per problem with thinking disabled, matching our greedy no-think probe;
the filter retains problems with ``room to improve,'' not guaranteed
failures. Baseline behavior is therefore consistent with the paper's own
calibration, and the freeze cannot be dismissed as a filter artifact.

\subsection{Real-data matching}

We used the ASSISTments 2009--2010 skill-builder data (401{,}756 rows,
4{,}217 students, 123 skills; officially free to use). Analysis units are
(user, skill) sequences with filters for main problems only, tutor mode, at
least five opportunities, and complete correctness records; the documented
order-identifier de-duplication was applied. The comparison pool is a seeded
sample of 5{,}000 eligible sequences. Domain mismatch is a stated limitation
(see Section~\ref{sec:prereg}): real trajectories span five opportunities on
different problems sharing a skill, while simulated trajectories span five
turns on one competition-mathematics problem. The comparison is
distributional, not item-matched.

\subsection{Output completeness}

Before any analysis, we verified output completeness: 300 of 300 episodes
present (6 $\times$ 50 exactly), five turns each, all probe vectors complete
and binary, 300 unique problems disjoint across archetypes, and all internal
consistency checks passing (rejected turns carry no student reply plus a
complaint; accepted turns carry replies; correctness vectors match baseline
plus probes). No partials and no gaps.

Preregistered exclusions: M2a excludes trajectories with no errors in the
pair window (12, 24, 13, 23, 12, 19 per archetype in Att./Sub./Ctr./Cau./Stp./Ver.
order; 1{,}146 of 5{,}000 real sequences); M3 excludes trajectories with
all-accepted or all-rejected turns (8, 3, 29, 21, 5, 6; 3{,}531 real). The KS
table in Appendix~A reports effective $n$ per test.

\section{Results}
\label{sec:results}

\subsection{The simulated students do not learn}
\label{sec:res-learn}

Figure~\ref{fig:curves} shows mean probe correctness by position. Real
students trace the classic learning curve, 0.499 $\pm$ 0.007 $\rightarrow$
0.688 $\pm$ 0.007 over five opportunities. Every simulated archetype is
flat: pooled correctness moves 0.347 $\rightarrow$ 0.343. In 297 of 300
trajectories, the probe vector is constant at its baseline value (195
all-wrong, 65.0\%; 102 all-correct, 34.0\%); the remaining three (1.0\%)
are mixed. The single
improving trajectory had zero accepted tutoring turns.
The two trajectories that got \emph{worse} after accepted tutoring
(\texttt{contrastive-comparison\_016}, \texttt{causal-justification\_044},
both correct at baseline) are the only hint in the data of tutoring-induced
harm, a thread for follow-up work rather than a claim.

\begin{figure}[t]
\centering
\includegraphics[width=\columnwidth]{figures/fig_learning_curves.png}
\caption{Mean correctness by probe (simulated) or opportunity (real). Real
students (ASSISTments, $n$=5{,}000) show the power-law learning curve; all
six simulated archetypes ($n$=50 each) are flat. Error bars show $\pm$SE.}
\label{fig:curves}
\end{figure}

\subsection{Degenerate distributional mismatch}
\label{sec:res-ks}

Table~\ref{tab:metrics} reports the preregistered metrics (mean $\pm$ SD).
The 24 KS tests are a preregistered formality: with the simulated
distributions collapsed to point masses, the p-values add nothing beyond
the descriptive degeneracy below, which is where the evidential weight
lies. The simulated distributions are degenerate: M1 slopes $\equiv$ 0.000
$\pm$ 0.000, M2a $\equiv$ 0.833 $\pm$ 0.000, and M3 $\equiv$ 0.000 $\pm$
0.000 in four archetypes, the other two showing only trace spread from
the three mixed trajectories. Real students show M1 mean slope $-0.35$
(SD 1.03), M2a 0.467 (SD 0.184), M2b mean streak 1.60 (SD 1.41), and M3
mean hint lift 0.087 (SD 0.498).

\begin{table}[t]
\centering
\caption{Preregistered metrics, mean $\pm$ SD. Simulated
$n$=50 per archetype; real $n$=5{,}000. Archetype abbreviations:
Att.=attempt-diagnosis, Sub.=subgoal-decomposition, Ctr.=contrastive-comparison,
Cau.=causal-justification, Stp.=step-demonstration, Ver.=independent-verification.
Post-exclusion counts are archived with the reproduction package.}
\label{tab:metrics}
\begin{tabular}{lcccc}
\toprule
Archetype & M1 slope & M2a persist. & M2b streak & M3 hint-resp. \\
\midrule
Att. & 0.000$\pm$0.000 & 0.833$\pm$0.000 & 3.80$\pm$2.14 & 0.000$\pm$0.000 \\
Sub. & 0.000$\pm$0.000 & 0.833$\pm$0.000 & 2.60$\pm$2.50 & 0.000$\pm$0.000 \\
Ctr. & 0.045$\pm$0.312 & 0.831$\pm$0.014 & 3.66$\pm$2.19 & $-$0.024$\pm$0.106 \\
Cau. & 0.035$\pm$0.244 & 0.832$\pm$0.006 & 2.68$\pm$2.48 & $-$0.011$\pm$0.061 \\
Stp. & 0.000$\pm$0.000 & 0.833$\pm$0.000 & 3.80$\pm$2.14 & 0.000$\pm$0.000 \\
Ver. & 0.000$\pm$0.000 & 0.833$\pm$0.000 & 3.10$\pm$2.43 & 0.000$\pm$0.000 \\
Real & $-$0.350$\pm$1.032 & 0.467$\pm$0.184 & 1.60$\pm$1.41 & 0.087$\pm$0.498 \\
\bottomrule
\end{tabular}
\end{table}

All 24 preregistered KS tests reject H0 after Holm correction (the full
table, with raw p-values, sample sizes, and bootstrap confidence intervals,
is archived with the reproduction package). Per the preregistered decision
rule, H1 is supported on all metric families---but the ``difference'' is
degenerate: the simulator does not behave like a \emph{different kind} of
student; it behaves like \emph{no student at all}.

\paragraph{Construct-validity note.} The preregistered comparison is
distributional, not construct-matched: simulated trajectories span five
turns on a single competition-mathematics problem, while real trajectories
span five opportunities across different problems sharing a skill; the M3
operationalizations likewise differ (accepted-vs-rejected guidance versus
hint-vs-no-hint, the latter confounded by help-seeking selection).
Rejection was therefore expected even for a competent single-problem
simulator. The evidential weight of this section lies not in the p-values
but in the descriptive degeneracy documented above: point-mass simulated
distributions where real students show spread and learning.

\subsection{Archetypes are behaviorally indistinguishable}
\label{sec:res-arch}

On M1 through M3 the six archetype distributions coincide (all
degenerate at the same values); a classifier given only these metrics could
not beat chance at telling archetypes apart. The differences that exist do not
constitute distinct learning dynamics: gate acceptance rates of 20--46\%
($\chi^2$=54.1, $p\approx2\times10^{-10}$, df=5; 95\% Wilson confidence intervals
archived; turns nested within episodes, so the $p$-value is descriptive) and
surface reply lengths of 811--1{,}070 characters (Kruskal--Wallis
$p\approx6\times10^{-7}$ with heavy overlap; likewise descriptive), and baseline accuracy of
0.24--0.48 (confounded, because problem sets are disjoint by preregistered
design). Sherpa's premise is that preference gates produce \emph{distinct
student profiles}; here they produce no learning dynamics at all.

\subsection{Exploratory: gate decomposition---half the rejections are judge failures}
\label{sec:res-gates}

Of 1{,}500 teacher turns, gate outcomes decompose as follows.

\begin{center}
\begin{tabular}{lcc}
\toprule
Outcome & Turns & Share \\
\midrule
Rejected by guidance gate & 258 & 17.2\% \\
Rejected by adaptive gate, judge parse failure (fail-closed) & 370 & 24.7\% \\
Rejected by adaptive gate, reasoned preference failure & 394 & 26.3\% \\
Accepted & 478 & 31.9\% \\
\bottomrule
\end{tabular}
\end{center}

\noindent\textbf{370 of 764 adaptive-gate rejections (48.4\%) are the 1.7B
judge's inability to emit valid tagged XML}, not preference judgments
(Figure~\ref{fig:gates}). Conditional on valid judge output,
the adaptive gate accepts 54.8\% of turns. A manual audit of 30 sampled
rejections (single rater; no inter-rater reliability measured, so the
reasoned-vs-parse split is approximate) found all 15
reasoned rejections defensible (for example, correctly flagging
single-alternative messages for the contrastive archetype); the other 15
were empty-reason parse failures. The
reasoned-reject rate (26.3\%) is the cleaner estimate of genuine gate
strictness---and it directly bounds how much of the headline 68.1\%
rejection rate survives the judge-downsizing caveat.

\begin{figure}[t]
\centering
\includegraphics[width=\columnwidth]{figures/fig_gate_decomp.png}
\caption{Gate-pipeline outcomes per archetype (250 turns each). Nearly half
(370 of 764, 48.4\%) of adaptive-gate rejections are fail-closed 1.7B-judge
parse failures (red), not preference judgments (blue).}
\label{fig:gates}
\end{figure}

\subsection{Exploratory: leaked boxed answers are all wrong; no dose--response}
\label{sec:res-boxed}

The fixed-strategy teacher was instructed never to reveal the final answer,
yet \textbf{141 of 1{,}500 (9.4\%) teacher messages contain boxed final
answers} (the final answer enclosed in a \texttt{\textbackslash boxed\{\}}
tag)---the 1.7B teacher emits them regardless. The guidance gate caught
80; 29 were adaptive-rejected; \textbf{32 were accepted into the student's
visible history}. In the 23 episodes containing an accepted boxed answer,
\textbf{zero probes improved}. Post-hoc checking of the 45 leaked boxed
answers inside the 32 accepted messages shows \textbf{all 45 were
numerically wrong} (0/45 correct against the ground truth): the 1.7B
teacher does not merely leak answers, it leaks \emph{wrong} answers (for
example, \texttt{causal-justification\_044} accepted $\boxed{11}$ for a
ground truth of $12$, then got worse). This null therefore indicts teacher
quality rather than testing whether tutoring content can move probes; the
clean test is the positive control in \S\ref{sec:res-control}. There is no
dose--response either: the correlation between accepted turns per episode
and probe improvement is $r=-0.04$ (95\% CI $[-0.15, 0.08]$, $n=300$;
Figure~\ref{fig:dose}); five accepted turns move the probe about as much
as zero (mean change 0.016, $-$0.009, $-$0.016, 0.000, 0.000, 0.000 for
zero through five accepted turns; $n$=62, 106, 63, 40, 19, 10).

\begin{figure}[t]
\centering
\includegraphics[width=\columnwidth]{figures/fig_dose_response.png}
\caption{Mean change in probe correctness by accepted tutoring turns per
episode. No dose--response ($r=-0.04$). Error bars show $\pm$SE.}
\label{fig:dose}
\end{figure}

\subsection{Exploratory: positive control---injected true answers move nothing}
\label{sec:res-control}

The preregistered protocol has no positive control: no configuration in
which probes were expected to move, so the freeze could in principle
reflect a broken probe loop rather than an unresponsive student. Post
hoc---unlike the teacher-leaked (and uniformly wrong) boxed answers in
\S\ref{sec:res-boxed}, here the injected answer is the ground-truth
one---we injected the ground-truth boxed answer as an accepted teacher
turn into the probe history of 162 of the 300 banked episodes and re-ran
the identical greedy probe plus two-stage grading
(\texttt{analysis/answer\_injection\_control.py}; logged in the deviation
log as a post-hoc addition). The control history is built from the
injected turn alone: it \emph{replaces}, not augments, the episode's
accepted tutoring history, so the injected answer faces no competing
history---the more generous test. Of the 106 episodes wrong at baseline,
\textbf{zero flipped to correct (0/106)}: per archetype, 0/38
attempt-diagnosis, 0/26 subgoal-decomposition, 0/36 contrastive-comparison,
0/6 causal-justification (95\% CI [0, 2.8\%]; rule of three, $3/106$). Of the 56 already correct, 55 stayed correct
and one flipped to wrong (1/56; 95\% CI [0.05\%, 9.5\%], consistent with
the probe-noise rate in \S\ref{sec:res-probe}). The run stopped at 162
episodes: a 2-CPU VM shared with other jobs ran each episode's probe plus
grading in roughly five minutes, and compute contention halted the run.
Episodes were processed in file order, so causal-justification stopped at
12 of 50 while step-demonstration and independent-verification (138
episodes combined) were never run.

Two caveats bound what this control shows. First, the probe prompt
instructs the student ``Now try to solve the problem from scratch:''---a
direction that plausibly tells the student to ignore the history above.
The 0/106 null is therefore consistent with an instruction-following
alternative (the student obeying ``from scratch'' rather than consulting
history) as well as with a genuine inability to use history content; this
probe design cannot separate the two, and no absolute ``channel is dead''
claim is warranted. Second, the control is partial: the 0/106 bound covers
the four archetypes run, not the two unrun ones. The probe loop itself is
not broken (it answers 34.7\% unaided from parametric knowledge), and the
control rules out weak tutoring as the sole explanation of the freeze for
the four archetypes tested---but the unresponsiveness claim is scoped to
this solve-from-scratch probe
instruction, not to student capacity in general. The freeze remains an
unresponsiveness of the frozen 1.7B student under this probe instruction, not evidence
about preference gating.

\subsection{Exploratory: greedy probes are not deterministic read-outs}
\label{sec:res-probe}

The preregistration assumed temperature-0.0 probes are deterministic state
read-outs. That assumption does not hold: among the 62 trajectories with
zero accepted turns (identical probe context all six times), probe
\emph{generations} varied in 54 (87\%) while graded \emph{verdicts}
flipped in only one. Generation-level stochasticity does not propagate to
verdicts, so the frozen-correctness finding is not an artifact of
re-scoring a single generation; the determinism assumption is dropped.

\section{Discussion}
\label{sec:discussion}

\paragraph{What failed.} Two candidate explanations fit the preregistered
data, which could not adjudicate between them. First, \emph{the tutoring
is vacuous}: with most teacher turns rejected and a weak 1.7B
fixed-strategy teacher (which leaks wrong boxed answers 9.4\% of the
time, all of them incorrect), the accepted histories that probes fork
from are thin (mean 1.59 accepted turns per episode; 62 episodes with
zero). We did not rate the pedagogical quality of the 478 accepted turns,
so teaching quality remains an open question. Second, \emph{the student
does not use its history}: the post-hoc positive control
(\S\ref{sec:res-control}) injects the true boxed answer as the probe
history (replacing, not augmenting, the accepted turns) and moves zero of
106 wrong-at-baseline probes (162 episodes, 95\% CI [0, 2.8\%]). This rules
out weak tutoring as the sole explanation, but it does not establish
general incapacity: the probe instruction ``Now try to solve the problem
from scratch'' plausibly tells the student to ignore history, so
instruction-following is an explicit alternative to incapacity, and the
control covers only four of six archetypes. The first hypothesis is true
descriptively; under the solve-from-scratch probe instruction, the student
is unresponsive to history content regardless of teaching quality. What the
data rules out is the middle position: that these archetype-conditioned
simulators, as implemented here, exhibit student-like learning dynamics.

\paragraph{Which findings survive which deviation.} The 24.7\% parse-failure
mass is the component most plausibly explained by the 1.7B judge---a 27B
judge should reduce it, but the paper itself reports that its 8B
training-time judge ``struggles to follow the required format in long
contexts'' \citep{xu2026sherpa}, so format-following failure is documented
at the paper's own judge scale and may be partly intrinsic rather than a
pure downsizing artifact. The 26.3\% reasoned-reject rate remains the
cleaner estimate of genuine gate strictness. The frozen probes, the null
dose--response, and the accepted-answer-reveal nulls do not depend on judge
quality. The fixed-strategy teacher is a deliberate scope choice as well as
a constraint: auditing the \emph{student} simulator requires holding
teaching fixed, because an adaptive teacher would confound student-side
validity with teacher quality. This is also the largest deviation from the
paper's own pipeline: its reported 20.5 percentage point improvement in
student test accuracy came from an RL-trained 8B teacher with a teacher
preparation stage and ten turns \citep{xu2026sherpa}; our fixed-strategy 1.7B
teacher, no preparation stage, and five turns deliberately hold none of that
constant. The post-hoc positive control (\S\ref{sec:res-control}) narrows
the two: bad teaching is real (thin accepted histories, wrong leaked
answers), but under this probe instruction the binding constraint is the
unresponsive student---even true answers injected into history move
nothing. The frozen-probe finding
therefore describes the student simulator under this scoped protocol, not
Sherpa's full system. The natural
follow-ups are a gates-disabled ablation, a pedagogical-quality rating of
accepted turns, and the 27B-judge replication.

\paragraph{What this means for simulated-student research.} The failure mode
we document is specific and actionable. Preference gating as a mechanism for
\emph{distinct student profiles} is untested by internal metrics: our
archetypes differ only in gate strictness and reply length, not in learning
behavior. Any evaluation of simulated students should include the cheapest
validity check available---a learning curve against a real-student
baseline---before more elaborate claims. And gate judges should be evaluated
on format reliability, not just judgment quality: nearly half of our adaptive
rejections were XML parse failures. The frozen-probe pattern is a concrete
instance of the competence paradox \citep{yuan2026}---the phenomenon in
which a capable model emulating a weaker learner performs at its own
capability floor, not the learner's: a broadly capable
model asked to emulate a partially knowledgeable learner answers at its
built-in capability level (34.7\% unaided) and cannot be moved off it by
instruction---the ``error pattern'' is a capability floor, not a student
profile.

\paragraph{The caveat, plainly.} The 1.7B downsizing is the sharpest threat
to our conclusions, and it cuts unevenly. It plausibly explains \emph{part}
of the gate dysfunction (the 24.7\% parse-failure mass should shrink under
the paper's 27B judge, though the paper documents format-struggle at its own
8B judge scale, so this mass may not dissolve entirely; the 26.3\% reasoned-reject mass is the cleaner
number). It cannot explain the frozen probes: the positive control shows the
student does not use even true history content under the solve-from-scratch
probe instruction---a scoped unresponsiveness of the weight-frozen 1.7B
student under this protocol (the control cannot separate
instruction-following from incapacity) rather than a gate effect. (The
accepted-answer-reveal null is confounded by the teacher leaking wrong
answers, \S\ref{sec:res-boxed}; the control supersedes it.) A
27B-judge replication is the obvious follow-up; we state exactly what we
expect it to change (gate acceptance up) and not to change (flat learning
curves).

\section{Limitations}
\label{sec:limitations}

Domain mismatch (competition mathematics versus K-12 skill builder;
preregistered); downsized 1.7B judge and teacher (the caveat above);
disjoint problem sets confound archetype baseline comparisons; probe
stochasticity is unmodeled (though verdict-stable); the language-model
answer judge's errors would propagate silently (no preregistered
spot-check); five turns versus the paper's ten; fifty trajectories per
archetype (CPU budget). A post-hoc positive control (injecting the true
boxed answer into the probe history; \S\ref{sec:res-control}) moved zero
of 106 wrong-at-baseline probes (162 episodes), so, under the
solve-from-scratch probe instruction, the freeze reflects the student's
unresponsiveness to history content---it cannot be attributed to
preference gating specifically, and the conclusions are scoped to this
frozen-student protocol.

\section{Conclusion}
\label{sec:conclusion}

We preregistered and executed an independent validity audit of Sherpa's
archetype-conditioned simulated students. The verdict is negative: 300
episodes, no measurable learning, archetypes indistinguishable on
learning-dynamics metrics, and a
gating pipeline in which half the rejections are judge format failures.
A post-hoc positive control (162 episodes, four of six archetypes; 0/106
flips, 95\% CI [0, 2.8\%]) shows the student does not use even true
answers under the solve-from-scratch probe instruction---though this
confounded control cannot separate instruction-following from incapacity.
This verdict is therefore scoped to the downsized frozen-student protocol:
it is a finding of unresponsiveness under this probe instruction, not an
attribution to preference gating and not a claim about Sherpa's full
system.
Negative results of this kind---preregistered, fully reproducible, with
every number traceable to banked artifacts---are how a research community
learns which of its instruments actually measure what they claim. The audit
code, data, preregistration, and one-command reproduction package are
publicly archived.

\section*{Declaration of Generative AI Software tools in the writing process}

During the preparation of this work, the author used Muse Spark (version
1.3, Meta) in all sections for drafting, analysis, and code validation,
and for grammar correction and proofreading throughout. After using this
tool, the author reviewed and edited the content as needed and takes full
responsibility for the content of the publication.

\section*{Competing Interests}

The author declares no competing interests. No funding was received for
this work.

\section*{Data and Code Availability}

The locked preregistration (\texttt{PREREGISTRATION.md}) with its
deviation log, the 300-episode trajectory file, the real-data comparison
pool, all analysis scripts, the post-hoc positive control, and a
one-command reproduction package (\texttt{bash reproduce.sh}) are archived
as supplementary material. A public code repository will be released upon
acceptance.

\begin{thebibliography}{99}
\bibitem[Xu et~al.(2026)]{xu2026sherpa}
XU, W., ZHANG, Y., WANG, Z.~Z., CHEN, C., AND YANG, D. 2026.
\newblock Sherpa: Teaching LLMs to teach adaptively.
\newblock \emph{arXiv preprint arXiv:2610.08778}.

\bibitem[Feng et~al.(2009)]{feng2009}
FENG, M., HEFFERNAN, N.~T., AND KOEDINGER, K.~R. 2009.
\newblock Addressing the assessment challenge with an online system that tutors
  as it assesses.
\newblock \emph{User Modeling and User-Adapted Interaction}, 19, 243--266.

\bibitem[Newell and Rosenbloom(1981)]{newell1981}
NEWELL, A. AND ROSENBLOOM, P.~S. 1981.
\newblock Mechanisms of skill acquisition and the law of practice.
\newblock In Anderson, J.~R., editor, \emph{Cognitive Skills and Their
  Acquisition}, pages 1--51. Lawrence Erlbaum Associates, Hillsdale, NJ.

\bibitem[Corbett and Anderson(1994)]{corbett1994}
CORBETT, A.~T. AND ANDERSON, J.~R. 1994.
\newblock Knowledge tracing: Modeling the acquisition of procedural knowledge.
\newblock \emph{User Modeling and User-Adapted Interaction}, 4, 253--278.

\bibitem[Nosek et~al.(2018)]{nosek2018}
NOSEK, B.~A., EBERSOLE, C.~R., DEHAVEN, A.~C., AND MELLOR, D.~T. 2018.
\newblock The preregistration revolution.
\newblock \emph{Proceedings of the National Academy of Sciences}, 115(11),
  2600--2606.

\bibitem[Scarlatos et~al.(2026)]{scarlatos2026}
SCARLATOS, A., LEE, J., WOODHEAD, S., AND LAN, A. 2026.
\newblock Simulated students in tutoring dialogues: Substance or illusion?
\newblock \emph{arXiv preprint arXiv:2601.04025}.

\bibitem[Lu and Wang(2024)]{lu2024}
LU, X. AND WANG, X. 2024.
\newblock Generative Students: Using LLM-simulated student profiles to support
  question item evaluation.
\newblock In \emph{Proceedings of the Eleventh ACM Conference on Learning at
  Scale (L@S '24)}. ACM.

\bibitem[Srivatsa et~al.(2025)]{srivatsa2025}
SRIVATSA, K.~V.~A., MAURYA, K.~K., AND KOCHMAR, E. 2025.
\newblock Can LLMs reliably simulate real students' abilities in mathematics
  and reading comprehension?
\newblock \emph{arXiv preprint arXiv:2507.08232}.

\bibitem[An et~al.(2026)]{an2026}
AN, Y., WANG, E., WANG, B., AND HASHMI, R. 2026.
\newblock From mastery profile to simulated response: Stochastic student
  knowledge graphs (SSKG) for faithful LLM student simulation.
\newblock \emph{arXiv preprint arXiv:2608.21668}.

\bibitem[Scarlatos et~al.(2025)]{scarlatos2025}
SCARLATOS, A., BAKER, R.~S., AND LAN, A. 2025.
\newblock Exploring knowledge tracing in tutor-student dialogues using LLMs.
\newblock In \emph{Proceedings of the 15th International Conference on
  Learning Analytics and Knowledge (LAK '25)}. ACM.

\bibitem[Yuan et~al.(2026)]{yuan2026}
YUAN, Z., XIAO, Y., LI, M., XUAN, W., TONG, R., DIAB, M., AND MITCHELL, T. 2026.
\newblock Towards valid student simulation with large language models.
\newblock \emph{arXiv preprint arXiv:2601.05473}.
\end{thebibliography}

\appendix

\section{Preregistered statistical tests: full results}

All 24 two-sample Kolmogorov--Smirnov tests (simulated $n$=50 per archetype,
minus preregistered exclusions, versus the 5{,}000-sequence real pool;
Holm--Bonferroni, $\alpha$=0.05) reject the null. The complete table with
test statistics, raw and adjusted p-values, effective sample sizes, and
bootstrap 95\% confidence intervals for median differences is archived in
\texttt{analysis/ks\_tests.json} in the reproduction package.

\section{Reproduction package contents}

The reproduction package archives: the locked preregistration
(\texttt{PREREGISTRATION.md}) with its deviation log; the 300-episode
trajectory file (\texttt{results/trajectories.jsonl}); the real-data
comparison pool and per-unit metrics; the analysis scripts that produced
every table and figure; the post-hoc answer-injection positive control
(\texttt{analysis/answer\_injection\_control.py} and
\texttt{analysis/answer\_injection\_control.jsonl}); the strengthening
analyses; and a verifier script
that re-checks the headline numbers. One command (\texttt{bash
reproduce.sh}) re-runs the full analysis from the banked artifacts.

\end{document}
